%%%%%%%%%%%%%%%%%%%%%%%%%%%%%%%%%%%%
% Slide options
%%%%%%%%%%%%%%%%%%%%%%%%%%%%%%%%%%%%

% Option 1: Slides with solutions

\documentclass[slidestop,compress,mathserif]{beamer}
\newcommand{\soln}[1]{\textit{#1}}
\newcommand{\solnGr}[1]{#1}

% Option 2: Handouts without solutions

%\documentclass[11pt,containsverbatim,handout]{beamer}
%\usepackage{pgfpages}
%\pgfpagesuselayout{4 on 1}[letterpaper,landscape,border shrink=5mm]
%\newcommand{\soln}[1]{ }
%\newcommand{\solnGr}{ }

%%%%%%%%%%%%%%%%%%%%%%%%%%%%%%%%%%%%
% Style
%%%%%%%%%%%%%%%%%%%%%%%%%%%%%%%%%%%%

\def\chp7@path{../../Chp 7}
\input{../../lec_style.tex}


%%%%%%%%%%%%%%%%%%%%%%%%%%%%%%%%%%%%
% Preamble
%%%%%%%%%%%%%%%%%%%%%%%%%%%%%%%%%%%%

\title[Lecture 24]{MA213: Lecture 24}
\subtitle{Module 4: Inference}
\author{OpenIntro Statistics, 4th Edition}
\institute{$\:$ \\ {\footnotesize Based on slides developed by Mine \c{C}etinkaya-Rundel of OpenIntro. \\
The slides may be copied, edited, and/or shared via the \webLink{http://creativecommons.org/licenses/by-sa/3.0/us/}{CC BY-SA license.} \\
Some images may be included under fair use guidelines (educational purposes).}}
\date{}


%%%%%%%%%%%%%%%%%%%%%%%%%%%%%%%%%%%%
% Begin document
%%%%%%%%%%%%%%%%%%%%%%%%%%%%%%%%%%%%

\begin{document}


%%%%%%%%%%%%%%%%%%%%%%%%%%%%%%%%%%%%
% Title page
%%%%%%%%%%%%%%%%%%%%%%%%%%%%%%%%%%%%

{
\addtocounter{framenumber}{-1} 
{\removepagenumbers 
\usebackgroundtemplate{\includegraphics[width=\paperwidth]{../../OpenIntro_Grid_4_3-01.jpg}}
\begin{frame}

\hfill 
\includegraphics[width=20mm]{../../oiLogo_highres}

\titlepage

\end{frame}
}
}

\include{Lecture24_agenda}
%%%%%%%%%%%%%%%%%%%%%%%%%%%%%%%%%%%%
% Sections
%%%%%%%%%%%%%%%%%%%%%%%%%%%%%%%%%%%%

\begin{frame}
\frametitle{Recall: CLT for the sample mean}

\formula{Central limit theorem (sample mean)}
{Sample means will be nearly normally distributed with mean equal to the population mean, $\mu$, and standard error equal to $\frac{\sigma}{\sqrt{n}}$.
\[ \bar{X} \sim N \pr{ mean = \mu, SE = \frac{\sigma}{\sqrt{n}} } \]
}

\begin{itemize}
    \item This is exactly true if the population distribution is Normal, regardless of sample size $n$ (we derived this in Lecture 16)
    % \pause
    \item But it is also true \emph{regardless of the shape of the population distribution}, as long as the sample size $n$ is sufficiently large.
    % \pause
    \item Different distributions will require different sample sizes for the CLT to ``kick in''.
    % \pause
    \item Rule of thumb is that $n \geq 30$ is sufficient for most distributions.
\end{itemize}

\end{frame}
%%%%%%%%%%%%%%%%%%%%%%%%%%%%%%%%%%%

\begin{frame}
\frametitle{Applying the CLT for the sample mean}

\dq{Which of the following scenarios would you be able to use the Normal distribution to approximate a 95\% confidence interval for the population mean $\mu$? In each case, you take a random sample of size $n$ from the population and calculate the sample mean $\bar{X}=49$.}
\begin{enumerate}[(a)]
\solnMult{Population: Normal ($\sigma = 20$). $n = 100$. }
\solnMult{Population: Normal ($\sigma = 20$). $n = 10$. }
  \item Population: Normal ($\sigma = $ unknown). $n = 10$.
\solnMult{Population: right skewed ($\sigma = 20$). $n = 100$. }
  \item Population: right skewed ($\sigma = 20$). $n = 10$. 
  \item Population: right skewed ($\sigma = $ unknown). $n = 10$.
\end{enumerate}
\end{frame}
%%%%%%%%%%%%%%%%%%%%%%%%%%%%%%%%%%%

\begin{frame}
\frametitle{Think/Pair/Share: Applying the CLT for the sample mean}

You take a random sample of size $n = 10$ and calculate the sample mean $\bar{X}=49$. Assume that the population is Normal with $\sigma = 20$. 

Calculate a 95\% confidence interval for the population mean $\mu$. Is the data consistent with a population mean of 60?

Hint: use \[ \bar{X} \sim N \pr{ mean = \mu, SE = \frac{\sigma}{\sqrt{n}} } \]

\begin{itemize}
    \item 2 minutes to think individually
    \item 2 minutes to discuss with a partner
    \item 2 minutes to share with the class
\end{itemize}

\end{frame}

%%%%%%%%%%%%%%%%%%%%%%%%%%%%%%%%%%%

\begin{frame}
\frametitle{Applying the CLT for the sample mean}
\dq{Assume that the conditions are met. If you know the population standard deviation $\sigma$, the sample size $n$, and the sample mean $\bar{X}$, how would you calculate a 95\% confidence interval for the population mean $\mu$?}

\soln{\only<2-3>{
\[point estimate \pm critical~value \times standard~error \]\\
\[\bar{X} \pm z^{\star} \times \frac{\sigma}{\sqrt{n}} \]
where $z^{\star}$ is the critical value from the standard Normal distribution for a 95\% confidence level (approximately 1.96).
}
}
\soln{\only<3>{
\[49 \pm 1.96 \times \frac{20}{\sqrt{10}} = 49 \pm 12.4= [36.6, 61.4]\]
}
}

\end{frame}

%%%%%%%%%%%%%%%%%%%%%%%%%%%%%%%%%%%
\begin{frame}
\frametitle{Applying the CLT for the sample mean}
\dq{Assume that the conditions are met. If you know the population standard deviation $\sigma$, the sample size $n$, and the sample mean $\bar{X}$, how would you design a hypothesis test at a significance level $\alpha$ to check whether the population mean is less than $\mu_0$?}

\soln{\only<2>{
$H_0: \mu = \mu_0$ \\
$H_A: \mu < \mu_0$ \\
Test statistic: \[ z_{obs} = \frac{\bar{X} - \mu_0}{\sigma / \sqrt{n}} \] \\
p-value: \[P(Z < z_{obs})\]\\
\texttt{pnorm(z\_obs)} \\
Check if p-value $< \alpha$ to make a conclusion.
}
}
\end{frame}

%%%%%%%%%%%%%%%%%%%%%%%%%%%%%%%%%%%
\begin{frame}
\frametitle{This Time}
What if the population is normal, but we don't know the population standard deviation, $\sigma$?

\begin{itemize}
    \item We can estimate $\sigma$ using the sample standard deviation $s$.
    \item However, this introduces additional uncertainty in our estimate of the standard error.
    \item If the sample size is large enough (Rule of thumb: $n \geq 30$), this additional uncertainty is negligible, because $s$ will be a good estimate of $\sigma$.
    \item But if the sample size is small, $s$ may not be a good estimate of $\sigma$.
    \item To account for this, we use a different distribution: the \hl{t-distribution}.
\end{itemize}

\end{frame}

%%%%%%%%%%%%%%%%%%%%%%%%%%%%%%%%%%%%


\section{One-sample means with the \texorpdfstring{$t$}{t} distribution}

%%%%%%%%%%%%%%%%%%%%%%%%%%%%%%%%%%%%


\begin{frame}
\frametitle{Guinness Quality Control: The Origin of the $t$-test}

\dq{William Sealy Gosset, working at the Guinness Brewery in the early 1900s, wanted to ensure the quality of stout. He would take a small sample of batches and measure a quality metric, (e.g. alcohol content). Suppose the target mean for alcohol content is 4.5\% ABV (percent alcohol by volume). Below is a sample of 8 batches.}

{\scriptsize
\begin{center}
\begin{tabular}{r|c}
Batch & Alcohol Content (\% ABV) \\
\hline
1 & 4.8 \\
2 & 4.6 \\
3 & 4.7 \\
4 & 4.5 \\
5 & 4.9 \\
6 & 4.6 \\
7 & 4.8 \\
8 & 4.4 \\
\end{tabular}
\end{center}
}

\vfill
\rule{2.5cm}{0.25pt} \\
{\tiny Inspired by the original work of W.S. Gosset ("Student") at Guinness Brewery.}

\end{frame}

%%%%%%%%%%%%%%%%%%%%%%%%%%%%%%%%%%%
\section{Edfinity Quiz: Setting up the Guinness test}
%%%%%%%%%%%%%%%%%%%%%%%%%%%%%%%%%%%


\begin{frame}
\frametitle{Guinness Quality Control: Framing the Question}

\begin{itemize}
\item We want to investigate if the mean alcohol content of Guinness stout batches differs from the target value of 4.5\% ABV.
\pause
\item One approach is to take a small sample of batches and compare the sample mean alcohol content to the target.
\pause
\item $H_0:$ The average alcohol content is equal to 4.5\% ABV. \\
$H_A:$ The average alcohol content is different from 4.5\% ABV.
\pause
\item test statistic: \[ \frac{\bar{X} - 4.5}{s / \sqrt{n}} \]
\end{itemize}

\orange{This test statistic looks similar to the $z$ test statistic, but uses $s$ instead of $\sigma$. How does this change things?}
% $\:$ \\

% \dq{Each row in the data set represents the alcohol content measured from a batch of Guiness stout. Are the batch measurements independent?}

% \soln{\pause Yes, we will assume that batches are produced independently.}

\end{frame}

%%%%%%%%%%%%%%%%%%%%%%%%%%%%%%%%%%%%
\section{R demo: What if we used a Normal model for \texorpdfstring{$\bar{X}$}{X-bar}?}
%%%%%%%%%%%%%%%%%%%%%%%%%%%%%%%%%%%

\begin{frame}
\frametitle{Review: what purpose does a large sample serve?}

As long as observations are independent, and the population distribution is not extremely skewed, a large sample would ensure that...

\begin{itemize}

\item the sampling distribution of the mean is nearly normal

\item the estimate of the standard error, as $\frac{s}{\sqrt{n}}$, is reliable

\end{itemize}

\pause
For small samples, we need to be more careful about both of these points.

\begin{itemize}
\item We need to verify that the \hl{population distribution is approximately normal.}
\begin{itemize}
\item Look at the data, check for skewness/outliers
\item Think about the context
\end{itemize}
\item If we don't know $\sigma$, we need to address the uncertainty in $s$.
\begin{itemize}
\item \hl{$t$ distribution}
\end{itemize}
\end{itemize}

\end{frame}

% \begin{frame}
% \frametitle{Hypotheses}

% \pq{What are the hypotheses for testing whether the mean alcohol content differs from the target value?}

% \begin{enumerate}[(a)]
% \item  \mathhl{H_0:} $\mu = 0$ \\
% \mathhl{H_A:} $\mu \ne 0$
% \solnMult{  \mathhl{H_0:} $\mu = 4.5$ \\
% \mathhl{H_A:} $\mu \ne 4.5$}
% \item  \mathhl{H_0:} $p = 4.5$ \\
% \mathhl{H_A:} $p \ne 4.5$
% \end{enumerate}

% \end{frame}

%%%%%%%%%%%%%%%%%%%%%%%%%%%%%%%%%%%

\begin{frame}
\frametitle{Conditions}
\twocol{0.6}{0.4}
{
\begin{itemize}
\item \hl{Independence:} We assume each batch is produced independently.
\pause
\item \hl{Normality ($n<30$):} $\:$ \\
\end{itemize}
}
{
    \vspace{0.5cm}
\includegraphics[width=\textwidth]{\chp7@path/7-1_one_t/figures/guiness/guinessHist}
}
\pause
{\small
\vspace{-0.25cm}
\begin{itemize}
\item The data doesn't appear skewed, but it's difficult to assess. 
\item For alcohol levels, we might expect scores to be fairly symmetric around the target.
\end{itemize}
}
\pause
\hl{We will assume that the data come from a Normal population}
\end{frame}

%%%%%%%%%%%%%%%%%%%%%%%%%%%%%%%%%%%


%%%%%%%%%%%%%%%%%%%%%%%%%%%%%%%%%%%

\subsection{The normality condition}

%%%%%%%%%%%%%%%%%%%%%%%%%%%%%%%%%%%

\begin{frame}
\frametitle{The normality condition}

We should exercise caution when verifying the normality condition for small samples. It is important to not only examine the data but also think about where the data come from. 
\begin{itemize}
\item For example, ask: would I expect this distribution to be symmetric, and am I confident that outliers are rare?
\end{itemize}

\end{frame}

%%%%%%%%%%%%%%%%%%%%%%%%%%%%%%%%%%%

\subsection{Introducing the \texorpdfstring{$t$}{t} distribution}

%%%%%%%%%%%%%%%%%%%%%%%%%%%%%%%%%%%

\begin{frame}
\frametitle{The $t$ distribution}

\begin{itemize}

\item When the population standard deviation is unknown (almost always) and the population is Normal, the uncertainty of the standard error estimate is addressed by using a new distribution: the \hl{$t$ distribution}.

\pause

\item This distribution also has a bell shape, but its tails are \hl{thicker} than the normal model's.

\pause

\item Therefore observations are more likely to fall beyond two SDs from the mean than under the normal distribution.

\pause

\item Extra thick tails are helpful for resolving our problem with a less reliable estimate the standard error (since $n$ is small)

\end{itemize}

\begin{center}
\includegraphics[width=0.5\textwidth]{\chp7@path/7-1_one_t/figures/tDistCompareToNormalDist/tDistCompareToNormalDist}
\end{center}

\end{frame}

%%%%%%%%%%%%%%%%%%%%%%%%%%%%%%%%%%%

\begin{frame}
\frametitle{The $t$ distribution (cont.)}

\begin{itemize}

\item Always centered at zero, like the standard normal ($z$) distribution.

\item Has a single parameter: \hl{degrees of freedom} (\mathhl{df}).

\end{itemize}

\begin{center}
\includegraphics[width=0.8\textwidth]{\chp7@path/7-1_one_t/figures/tDistConvergeToNormalDist/tDistConvergeToNormalDist}
\end{center}

\pause

\dq{What happens to shape of the $t$ distribution as $df$ increases?}

\soln{\pause Approaches normal.}

\end{frame}

%%%%%%%%%%%%%%%%%%%%%%%%%%%%%%%%%%%

\subsection{Evaluating hypotheses using the \texorpdfstring{$t$}{t} distribution}

%%%%%%%%%%%%%%%%%%%%%%%%%%%%%%%%%%%


\section{R Demonstration: Try with \texorpdfstring{$t$}{t} distribution}

%%%%%%%%%%%%%%%%%%%%%%%%%%%%%%%%%%%

\begin{frame}
\frametitle{Back to Guinness: Sample Statistics}

\begin{itemize}
\item Sample mean ($\bar{x}$): 4.66
\item Sample standard deviation ($s$): 0.17
\item Sample size ($n$): 8
\end{itemize}

\begin{align*}
\orange{$\bar{x} = 4.66$} \\
\orange{s = 0.17} \\
\orange{n = 8}
\end{align*}

\end{frame}

%%%%%%%%%%%%%%%%%%%%%%%%%%%%%%%%%%%


\begin{frame}
\frametitle{Finding the test statistic}

\formula{Test statistic for inference on a small sample mean}
{The test statistic for inference on a small sample ($n < 30$) mean is the $t$ statistic with $df = n - 1$.
\[ t_{obs} = \frac{\text{point estimate} - \text{null value}}{SE} \]}

\pause

\vspace{-0.5cm}

\hl{in context...}
\begin{eqnarray*}
point~estimate &=& \bar{x} = 4.66 \\
\pause
SE &=& \frac{s}{\sqrt{n}} = \frac{0.17}{\sqrt{8}} = 0.060 \\
\pause
t_{obs} &=& \frac{4.66 - 4.5}{0.060} = 2.73 \\
\pause
df &=& 8 - 1 = 7
\end{eqnarray*}

\vspace{-0.25cm}

\Note{Null value is 4.5 because in the null hypothesis we set $\mu = 4.5$.}

\end{frame}

%%%%%%%%%%%%%%%%%%%%%%%%%%%%%%%%%%%

\begin{frame}[fragile]
\frametitle{Finding the p-value}

\begin{itemize}
\item The p-value is calculated as the tail area under the $t$ distribution.
\pause
\item Using R:
\begin{verbatim}
> 2 * pt(2.73, df = 7, lower.tail = FALSE)
[1] 0.029
\end{verbatim}
\pause
\item Visualized below:
\begin{center}
\includegraphics[width=0.7\textwidth]{\chp7@path/7-1_one_t/figures/guiness/guinessPvalue}
\end{center}
\pause
\item Or when these aren't available, we can use a $t$-table.
\end{itemize}
\end{frame}

%%%%%%%%%%%%%%%%%%%%%%%%%%%%%%%%%%%

\section{R Demonstration: Try with \texorpdfstring{$t$}{t} distribution}

%%%%%%%%%%%%%%%%%%%%%%%%%%%%%%%%%%%

\begin{frame}
\frametitle{Conclusion of the test}

\dq{What is the conclusion of this hypothesis test?}

\pause

Since the p-value is low, we have evidence that the mean alcohol content differs from the target value of 4.5\% ABV.

\end{frame}

%%%%%%%%%%%%%%%%%%%%%%%%%%%%%%%%%%%%

%%%%%%%%%%%%%%%%%%%%%%%%%%%%%%%%%%%

%%%%%%%%%%%%%%%%%%%%%%%%%%%%%%%%%%%%
% End document
%%%%%%%%%%%%%%%%%%%%%%%%%%%%%%%%%%%%

\end{document}